../cictikz/src/cictikz/data/tex/cictikz_preamble.tex

% From OTA to op amp: add a class AB output stage so the amplifier can
% drive a resistive load. The AB block level shifts the gate drives so
% both output devices conduct a small quiescent current.

\begin{circuitikz}[american, thick, transform shape, circuit ee IEC]

  ../cictikz/src/cictikz/data/tex/cictikz_lib.tex

  % the OTA as a triangle
  \draw (0,1.1) -- (0,-1.1) -- (2.0,0) -- cycle;
  \draw (-1.0,0.55) node[anchor=east] {$v_p$} to [short,o-] (0,0.55);
  \draw (-1.0,-0.55) node[anchor=east] {$v_n$} to [short,o-] (0,-0.55);
  \node[anchor=west, scale=0.9] at (0.15,0) {OTA};

  % class AB level shift block
  \draw (2.0,0) -- (2.6,0);
  \draw (2.6,-0.7) rectangle (3.6,0.7);
  \node[anchor=center, scale=0.9] at (3.1,0) {AB};

  % output devices: PMOS common source up, NMOS common source down
  \draw (5.0,0) \lvpmos{MO1}{};
  \draw (5.0,\grid) \vsupply;
  \draw (5.0,-\grid) \lvnmos{MO2}{};
  \draw (5.0,-\grid) \vground;

  \coordinate (gp) at (MO1.G);
  \coordinate (gn) at (MO2.G);
  \draw (3.6,0.45) -- (4.0,0.45) -- (4.0,0.8) -- (gp);
  \draw (3.6,-0.45) -- (4.0,-0.45) -- (4.0,-0.8) -- (gn);

  % output and the resistive load
  \fill (5.0,0) circle (0.075);
  \draw (5.0,0) -- (6.4,0);
  \draw (6.4,0) to [short,*-o] ++(0.0,0.0);
  \node[anchor=south] at (6.4,0.15) {$v_o$};
  \draw (6.4,0) -- (7.4,0);
  \draw (7.4,-1.6) \vresistor{$R_L$};
  \draw (7.4,-1.6) \vground;

\end{circuitikz}
\end{document}
